% Chapter 1: Introduction

\chapter{Introduction}

\section{Background and Motivation}

The last three years have been a genuine revolution in AI, led by the rapid adoption of large language models. Models from OpenAI, Gemini, Anthropic and other players have shown remarkable capabilities in language understanding, generation and reasoning that have pushed organisations to rethink how they approach knowledge-intensive work.
One of the most important architectural developments from this era is Retrieval-Augmented Generation (RAG). It was first formalised by \citet{lewis2020retrieval} to address two key limitations of standalone LLMs: hallucinations and the inability to access knowledge beyond their training data. By pairing a retriever with a generative model, RAG systems can ground their outputs in external databases, access up-to-date domain information, and have become the go-to architecture for AI applications in sectors like finance, legal and healthcare. In the beginning, RAG seemed like a practical solution: give the model the right documents, and it should produce a more grounded and reliable answer.
But with RAG came an evaluation problem. How do organisations actually verify their RAG systems are working as expected? Traditional methods, such as humans annotating outputs one by one or reference-based scores like BLEU and ROUGE, are either too expensive at scale or too blunt to capture the nuances of generative outputs. This gap gave rise to the LLM-as-a-Judge concept, where an LLM is used to score the outputs of another model across different quality criteria. Frameworks like RAGAS \citep{es2024ragas} and G-Eval \citep{liu2023geval} are now widely used for this kind of automated evaluation.
LLM judges are attractive because they scale, they are consistent, and they roughly align with human evaluators at a fraction of the cost. But this whole approach rests on one assumption: that the context the judge is given to evaluate against is trustworthy. In reality, RAG systems are vulnerable to corpus poisoning, where adversarial passages are injected into the knowledge base to steer the system toward wrong answers. When a judge evaluates a response against poisoned context, it may score it as faithful because the response does match the context, even though the context itself is fabricated.
Recent research has begun to expose the fragility of this assumption. Studies by \citet{raina2024llm} showcased that short adversarial phrases can systematically fool LLM judges into assigning higher scores, regardless of the actual quality of the text. \citet{zou2025poisonedrag} introduced the PoisonedRAG framework, highlighting that injecting small malicious pieces of text into a RAG knowledge base can manipulate RAG behaviour with an ASR (Attack Success Rate) of 90\% across different benchmarks. Further study by \citet{eiras2025know} showed that commonly used LLM judges are sensitive to distribution patterns and adversarial modifications in output. TrustRAG \citep{zhou2025trustrag} and RAGuard \citep{cheng2025raguard} proposed defence strategies, but they focused on generation and retrieval quality rather than evaluation reliability. Altogether, research and these findings point to a vulnerability in the system that questions the trustworthiness of LLM judges. Whether the judge itself is compromised by corpus poisoning and whether upstream filtering can restore it has not been studied systematically.

\subsection{The Scale of the Evaluation Challenge}

To better understand the urgency of this problem, let's assess the scale at which RAG evaluation is being conducted today across the industry. A typical enterprise RAG system handles thousands of queries every single day, going through massive databases with millions of documents to find related and relevant data and generate output. Manual evaluation of even a small sample of that will require a team of experts going through each response and annotating it; that's costly and too slow a process. To keep up, the industry-wide adoption of the LLM-as-a-Judge strategy began. But there is a catch: the security around these judges is still a black box.
The consequences of evaluation failure are not just academic. If these automated judges get it wrong, the real-world fallout can be huge. In finance, it could mean heavy regulatory fines. In healthcare, it might result in medically dangerous advice slipping through the cracks. In the legal world, it could lead to flat-out wrong guidance.
This creates a dangerous double blind spot. If an attacker manages to sneak bad information into the system, the AI judge might actually greenlight the wrong answer, reporting that everything is fine. Because standard safety checks fail to catch this, attackers are actually encouraged to poison the data. They know the compromised judge will cover their tracks while the system looks perfectly healthy.

\subsection{Origins of the Research}

The idea for this research came directly from my own time working in the industry. During my internship with the AIR Tech team at BNP Paribas, I was assigned to explore LLM as a Judge concept and build an automated AI-judge pipeline for compliance use cases. It was easy to see why companies prefer the efficiency of this setup, but I also saw firsthand just how fragile it really was.
Two big realisations drove this project. The first was that in highly regulated fields like banking, the cost of a missed error is massive. If an AI judge accidentally stamps a factually wrong answer as ``correct,'' that error gets baked into official audit trails, giving everyone a false sense of security.
The second realisation is that the AI judge has a fundamental flaw: it cannot double-check the facts it is handed. It only checks if the generated answer matches the provided documents. If an attacker has poisoned those documents, the judge is completely fooled, and no amount of clever prompt engineering can fix that.
Realising this major vulnerability that academics had not fully explored yet made it clear that we need practical, real-world defences that standard engineering teams can actually implement.

\section{Problem Statement}

This thesis aims to answer the following basic question: Does a conventional LLM judge used within a RAG evaluation pipeline still hold up against corpus poisoning, and to what extent will lightweight upstream filtering defend against corpus poisoning?

In a RAG evaluation pipeline, an LLM judge is used to rate a generated answer with regard to the retrieved context, for example evaluating whether the answer is faithful to the retrieved evidence during evaluation. Under adversarial conditions, the retrieved evidence may have been poisoned and may contain fake or harmful content. An answer is then regarded as supported by the retrieved context even though the evidence may contain a lot of irrelevant but factually supported but fabricated or maliciously created context. Since an LLM judge has no access to an external source of ground truth, it will only score an answer based on the provided context.
A straightforward mitigation is applying some kind of filter to retrieve context before feeding to the judge. Filter mechanisms like statistical anomaly detection, embedding based filter, adversarial classifier and LLM based moderation were all proposed and investigated. Yet whether they can actually improve downstream judge robustness when the corpus poisoning is coherent and aligned with the retriever is still unclear.
Also, many of the defence strategies proposed for RAG security such as training safety models, model fine-tuning, creating adversarial retriever or using ensemble defence mechanisms, might be computationally expensive, need dedicated datasets and require significant engineering expertise, making them difficult to be implemented by a single researcher, a small team, or on resource constrained machines. Efficient and lightweight defences, and empirical study of these defences, is thus highly needed.

\section{Research Questions}

The four research questions below take the thesis from diagnosis to mitigation design to analysis.

\textbf{Research Question 1 (RQ1)- Baseline Vulnerability:} How well do present LLM judges (GPT-5.4-nano, Gemma-4-26B-A4B-IT, DeepSeek-V3.2) score poisoned RAG triplets on RAGAS dimensions, and how does this vulnerability vary for different attack types?

\textbf{Research Question 2 (RQ2)- Comparative Judge Behaviour:} In what respects are the three judges different in terms of vulnerability profile, score distributions, and calibration behaviour under corpus poisoning? What does this tell us about the structural and training aspects that impact a model's behaviour as a judge?

\textbf{Research Question 3 (RQ3)- Pre-filter Efficacy:} Are light-weight upstream pre-filter pipelines (statistical multi-signal, or LLM-based) likely to help judges to be more reliable with no (or little) degradation on clean context performance?

\textbf{Research Question 4 (RQ4)- Limits of Filtering:} What effect do the different filtering mechanisms have on the rest of the context supplied to down-stream judges, and what trade-offs are made when they are applied?

\section{Contributions}

This thesis makes four principal contributions to the fields of LLM evaluation, RAG security, and automated quality assessment.

\textbf{First, the Poisoned Evaluation Dataset (v2).} A dataset of 2,400 question-context-answer triplets extracted from HotpotQA with different poisoning methods and at varying degrees of corruption. Specifically created to examine LLM judge behaviour in the presence of an adversarially manipulated evaluation corpus. We also provide the poisoning strategy to easily reproduce the same effect on other datasets.

\textbf{Second, a multi-metric meta-evaluation framework.} A framework to examine judge robustness to corpus poisoning, used to evaluate 3 modern LLM judges on conventional RAGAS metrics while also characterising their calibration behaviour, score distributions, and susceptibility to attack-specific behaviours.

\textbf{Third, a lightweight statistical pre-filter.} A lightweight, statistically-based pre-filter composed of several engineered features combined via an XGBoost classifier, evaluated alongside an LLM-based filter to study individual filtering performance and impact on downstream judges.

\textbf{Fourth, an end-to-end residual-aware evaluation framework.} An end-to-end framework inspired by the True/Survived decomposition, designed to determine whether a filter actually reduces the impact of a poison attack, as opposed to simply changing perceived faithfulness by removing distracting material.

Together with a comparative discussion of existing RAG defence approaches, these contributions advance the study of reliable automated evaluation under adversarial retrieval conditions.

\section{Thesis Structure}

The structure of the remainder of this thesis is organised as follows: Chapter 2 introduces the related work of RAG evaluation, corpus poisoning and the LLM-as-a-Judge framework, in which we identify the research gap this thesis is concerned with. Chapter 3 describes the proposed methods, including the dataset construction, the baseline meta-evaluation frame, the designs of pre-filters, and the validation setup. Chapter 4 provides the experimental results and their analysis. Chapter 5 discusses the impact of the findings, the limitations of the proposed methods and some implications of the study for trustworthy RAG evaluation. Chapter 6 concludes the thesis and proposes the directions of future work.
